\documentclass[10pt]{article}
\usepackage[margin=0.72in]{geometry}
\usepackage{amsmath}
\usepackage{amssymb}
\usepackage{booktabs}
\usepackage{graphicx}
\usepackage{caption}
\usepackage{microtype}
\title{SAC versus action-conditioned JEPA on the fractal terrain task}
\author{}
\date{RTX 3090 runs, September 20, 2026}
\begin{document}
\maketitle

\section*{Question and result}
All runs use the same centered-Gaussian squared-potential terrain, random
starts, $P_{\max}=6$ longitudinal power cap, 16 environments, held-out terrain
seeds, and 10,000-transition budget. SAC alone is the strongest observed
controller. The normalized, action-conditioned JEPA branch is mechanically
valid---it beats shuffled-action and identity diagnostics---but still reduces
held-out speed at this budget.

\begin{center}
\small
\begin{tabular}{p{1.6in}ccp{2.0in}}
\toprule
Run & Held-out speed at 10k & Difference from SAC & Status \\
\midrule
SAC-only & $1.655\pm0.056$ & --- & baseline \\
SAC+JEPA, old MSE target & $1.278\pm0.042$ & $-22.8\%$ & superseded cadence/scale safeguards \\
SAC+JEPA, corrected MSE & $1.403\pm0.020$ & $-15.2\%$ & variance floor and slow EMA \\
SAC+JEPA, cosine predictive loss & $1.188$ & $-28.2\%$ & scale-normalized diagnostic run \\
\bottomrule
\end{tabular}
\end{center}

\section*{Full corrected SAC+JEPA loss}
For a replayed transition $(o_t,a_t,o_{t+1})$, the actor has online encoder
$z_t=f_\theta(o_t)$, EMA target $\bar z_{t+1}=\bar f(o_{t+1})$, and predictor
$q_\psi$. The optimized actor objective is
\begin{equation}
\begin{aligned}
\mathcal L_{\rm actor}={}&\mathbb E\!\left[\alpha\log\pi(a|o)-\min(Q_1(o,a),Q_2(o,a))\right]\\
&+\lambda\left(1-\cos\!\left(q_\psi([z_t,a_t]),\operatorname{sg}[\bar z_{t+1}]\right)\right)\\
&+c_{\rm var}\frac{1}{d}\sum_{j=1}^d\max\!\left(0,s_{\rm min}-\operatorname{std}_B(z_{:,j})\right).
\end{aligned}
\end{equation}
The cosine loss is scale-free; the variance floor prevents a zero-variance
representation. The final run uses $\lambda=0.025$, $c_{\rm var}=0.01$,
$s_{\rm min}=0.1$, and updates the EMA target by 0.01 every 64 collected
transitions. Logged diagnostics include the weighted JEPA and variance terms,
action-shuffled loss, identity baseline, and embedding spread.

\begin{figure}[p]
\centering
\includegraphics[width=.49\linewidth]{../output/pdf/sac_only_learning_curve.png}
\includegraphics[width=.49\linewidth]{../dumps/sac_jepa_cosine_10k/plots/learning_curve.png}
\caption{Performance curves. Left: SAC-only. Right: normalized cosine SAC+JEPA.
The right lower panel is the scale-free predictive loss.}
\end{figure}

\begin{figure}[p]
\centering
\includegraphics[width=.49\linewidth]{../dumps/sac_10k/plots/latent_umap_sac_000010000_seed100001.png}
\includegraphics[width=.49\linewidth]{../dumps/sac_jepa_cosine_10k/plots/latent_umap_sac_000010000_seed100001.png}
\caption{Endpoint terrain-overlay trajectories and UMAPs on the same held-out
seed. Left: SAC-only. Right: normalized cosine SAC+JEPA.}
\end{figure}

\begin{figure}[p]
\centering
\includegraphics[width=.49\linewidth]{../dumps/sac_10k/plots/dynamics_history_sac_000010000_seed100001.png}
\includegraphics[width=.49\linewidth]{../dumps/sac_jepa_cosine_10k/plots/dynamics_history_sac_000010000_seed100001.png}
\caption{Full velocity, control, and acceleration diagnostics. Left: SAC-only.
Right: normalized cosine SAC+JEPA.}
\end{figure}

\begin{figure}[p]
\centering
\includegraphics[width=.49\linewidth]{../dumps/sac_10k/plots/encoder_cka_matrix_seed100001.png}
\includegraphics[width=.49\linewidth]{../dumps/sac_jepa_cosine_10k/plots/encoder_cka_matrix_seed100001.png}
\caption{Encoder evolution across the five saved checkpoints, measured by
linear CKA on a common held-out inference trajectory.}
\end{figure}

\begin{figure}[p]
\centering
\includegraphics[width=.49\linewidth]{../dumps/sac_10k/plots/pairwise_encoder_umap_5x5_seed100001.png}
\includegraphics[width=.49\linewidth]{../dumps/sac_jepa_cosine_10k/plots/pairwise_encoder_umap_5x5_seed100001.png}
\caption{Checkpoint-pair latent matrix on the shared deterministic held-out
trajectory. Each cell jointly embeds the row and counterpart encoder states
(blue/orange); diagonal cells are self-comparisons. Left: SAC-only. Right:
normalized cosine SAC+JEPA.}
\end{figure}

\section*{Interpretation}
The cosine replacement fixes the earlier loss-scale failure: the predictive
term is no longer numerically negligible relative to the variance term, and
the shuffled-action loss exceeds the real-action loss ($0.0099$ versus $0.0038$
at 10k). Thus the predictor uses actions. Nonetheless, performance is lower
than SAC alone. At this horizon and budget, the evidence favors SAC-only;
multi-step targets and matched multi-seed sweeps remain the next JEPA test.

\clearpage
\section*{V2: five-times-longer SAC+JEPA run (50k transitions)}
V2 repeats the corrected cosine SAC+JEPA configuration from V1 on the RTX
3090, with the same terrain split, power limit, random starts, and held-out
evaluation seeds. Only the transition budget changes: 10k to 50k. Checkpoints
are saved every 10k, so all representation plots compare five within-run
checkpoints on the same held-out deterministic trajectory.

\begin{center}
\small
\begin{tabular}{lcc}
\toprule
Run & Held-out speed & Interpretation \\
\midrule
V1 cosine SAC+JEPA, 10k & $1.188$ & early-budget endpoint \\
V2 cosine SAC+JEPA, 50k & $2.954\pm0.049$ & continued improvement, plateau after about 18k \\
Matched SAC-only, 50k & running & required before asserting relative slowness \\
\bottomrule
\end{tabular}
\end{center}

The auxiliary branch remains active in V2: at 10,240 transitions its real
action-conditioned cosine loss is $0.0045$ versus $0.0101$ with actions
shuffled. The gap narrows late in training, as the controller and observations
become more predictable. Thus V2 establishes that the JEPA controller keeps
learning with a larger budget, but it does \emph{not} yet prove that it is
slower than SAC: that requires the matched 50k SAC-only control now in flight.

\begin{figure}[p]
\centering
\includegraphics[width=.92\linewidth]{../dumps/sac_jepa_cosine_50k_v2/plots/learning_curve.png}
\caption{V2 held-out performance and auxiliary diagnostics across 50k
transitions. The held-out speed rises sharply after 10k and stabilizes near
2.95 from approximately 18k onward.}
\end{figure}

\begin{figure}[p]
\centering
\includegraphics[width=.49\linewidth]{../dumps/sac_jepa_cosine_50k_v2/plots/latent_umap_sac_000050000_seed100001.png}
\includegraphics[width=.49\linewidth]{../dumps/sac_jepa_cosine_50k_v2/plots/dynamics_history_sac_000050000_seed100001.png}
\caption{V2 endpoint deterministic inference on held-out seed 100001. Left:
terrain-overlay trajectory and encoder UMAP. Right: velocity, input, field,
damping, and net acceleration histories.}
\end{figure}

\begin{figure}[p]
\centering
\includegraphics[width=.49\linewidth]{../dumps/sac_jepa_cosine_50k_v2/plots/encoder_cka_matrix_seed100001.png}
\includegraphics[width=.49\linewidth]{../dumps/sac_jepa_cosine_50k_v2/plots/pairwise_encoder_umap_5x5_seed100001.png}
\caption{V2 encoder evolution across the five 10k-spaced checkpoints on one
shared held-out trajectory. Left: linear CKA. Right: pairwise joint UMAP
matrix (blue/orange are the two checkpoint encoders in each cell).}
\end{figure}

\clearpage
\section*{V2 completion: matched 50k SAC-only control}
The control mentioned above has now completed under the same 50k transition
budget, terrain split, power limit, random-start protocol, evaluation seeds,
and checkpoint schedule. This is the valid comparison for the question
``does JEPA merely learn more slowly?''

\begin{center}
\small
\begin{tabular}{lccc}
\toprule
Transitions & SAC-only & SAC+JEPA & JEPA minus SAC \\
\midrule
10,240 & $1.426$ & $1.402$ & $-0.024$ \\
12,288 & $2.258$ & $2.433$ & $+0.175$ \\
16,384 & $2.836$ & $2.844$ & $+0.008$ \\
50,000 & $2.952\pm0.030$ & $2.954\pm0.049$ & $+0.003$ \\
\bottomrule
\end{tabular}
\end{center}

\begin{figure}[h]
\centering
\includegraphics[width=.88\linewidth]{../dumps/sac_jepa_cosine_50k_v2/plots/matched_sac_vs_jepa_50k.png}
\caption{Direct matched 50k held-out comparison. The curves cross during the
early transient and are nearly identical from about 16k onward. Shading is the
held-out standard deviation across the fixed evaluation seeds.}
\end{figure}

\noindent\textbf{Conclusion.} In this one matched seed/configuration pair,
the corrected cosine SAC+JEPA branch is neither persistently worse nor
detectably slower than SAC-only. Its final mean differs by $0.003$ speed
($0.09\%$), far smaller than the reported evaluation spread. This establishes
convergence at 50k; it does not establish a general JEPA benefit. Matched
multi-seed replicates are required for that claim.

\clearpage
\section*{V2 visual companion: SAC-only versus SAC+JEPA}
These panels are the corresponding standard inference artifacts from the two
matched 50k endpoints. Both use the same held-out terrain seed (100001) and
the same deterministic-inference protocol; only the learned controller and
encoder differ.

\begin{figure}[p]
\centering
\includegraphics[width=.49\linewidth]{../dumps/sac_50k_control_v2/plots/latent_umap_sac_000050000_seed100001.png}
\includegraphics[width=.49\linewidth]{../dumps/sac_jepa_cosine_50k_v2/plots/latent_umap_sac_000050000_seed100001.png}
\caption{Matched endpoint terrain-overlay trajectory and latent UMAP. Left:
SAC-only. Right: SAC+JEPA. Point colour maps the two UMAP coordinates back to
the physical trajectory.}
\end{figure}

\begin{figure}[p]
\centering
\includegraphics[width=.49\linewidth]{../dumps/sac_50k_control_v2/plots/dynamics_history_sac_000050000_seed100001.png}
\includegraphics[width=.49\linewidth]{../dumps/sac_jepa_cosine_50k_v2/plots/dynamics_history_sac_000050000_seed100001.png}
\caption{Matched endpoint velocity and acceleration histories. Left:
SAC-only. Right: SAC+JEPA. Each panel contains speed and latent velocity,
actions, and signed input/potential/damping/net acceleration components.}
\end{figure}

\begin{figure}[p]
\centering
\includegraphics[width=.49\linewidth]{../dumps/sac_50k_control_v2/plots/encoder_cka_matrix_seed100001.png}
\includegraphics[width=.49\linewidth]{../dumps/sac_jepa_cosine_50k_v2/plots/encoder_cka_matrix_seed100001.png}
\caption{Within-run representation evolution on a common held-out trajectory.
Left: SAC-only. Right: SAC+JEPA. Each 5 by 5 matrix is linear CKA across the
10k, 20k, 30k, 40k, and 50k checkpoints.}
\end{figure}

\begin{figure}[p]
\centering
\includegraphics[width=.49\linewidth]{../dumps/sac_50k_control_v2/plots/pairwise_encoder_umap_5x5_seed100001.png}
\includegraphics[width=.49\linewidth]{../dumps/sac_jepa_cosine_50k_v2/plots/pairwise_encoder_umap_5x5_seed100001.png}
\caption{Pairwise joint UMAP trajectory matrices for the same checkpoints.
Blue and orange are the row and counterpart encoders evaluated on identical
held-out inference states. Left: SAC-only. Right: SAC+JEPA.}
\end{figure}

\clearpage
\section*{Post-run diagnostic fidelity correction}
The original V2 run is unchanged. Its speed curve, endpoint evaluations, and
inference/representation artifacts remain valid. Two reported auxiliary
details are corrected for future runs: (i) the cosine-prediction panel is now
labelled ``cosine distance ($1-\cos$)'', not MSE; and (ii) training now logs
the mean of every replay update since the prior evaluation, rather than the
final minibatch alone. Consequently, the V2 single-checkpoint real-action
versus shuffled-action loss values are diagnostic snapshots, not evidence of
a statistically significant action effect. The corrected identity baseline
also compares target-encoder embeddings of $o_t$ and $o_{t+1}$, isolating
temporal predictability from EMA lag. These corrections do not change the
matched 50k performance conclusion; they improve fidelity for subsequent
JEPA ablations.

\clearpage
\section*{V3: boundary-safe eight-step JEPA with inverse dynamics}
V3 changes the auxiliary task, not the SAC environment or reward. Replay now
marks every episode boundary (including time-limit truncations), so an
eight-action window cannot splice unrelated reset terrains. The predictor
receives $[z_t,a_t,\ldots,a_{t+7}]$ and targets the EMA encoding of
$o_{t+8}$. An inverse head predicts $a_t$ from the two latent endpoints. The
optimized auxiliary terms are
\[
\lambda\,[1-\cos(q([z_t,a_{t:t+7}]),\operatorname{sg}[\bar z_{t+8}])]
+\lambda_{\rm inv}\,\|g(z_t,\bar z_{t+8})-a_t\|_2^2,
\]
with $\lambda=\lambda_{\rm inv}=0.025$. V3 also logs the interval-mean
real/shuffled ratio, rather than a final-minibatch snapshot.

\begin{center}\small
\begin{tabular}{lcc}
\toprule
Run & Held-out speed at 10k & Real/shuffled ratio \\
\midrule
V2 cosine SAC+JEPA & $1.188$ & late-run ratio approached 1 \\
V3 K=8 + inverse SAC+JEPA & $0.863$ & $0.204$ \\
\bottomrule
\end{tabular}\end{center}

V3 clears the action-conditioning acceptance test: the shuffled predictor is
substantially worse than the real-action predictor at 10k. It does not yet
improve early control performance, so this result isolates a repaired
representation objective rather than a performance win.

\begin{figure}[p]\centering
\includegraphics[width=.88\linewidth]{../dumps/sac_jepa_v3_10k_corrected/plots/learning_curve.png}
\caption{V3 learning curve and auxiliary diagnostic over 10k transitions.}
\end{figure}

\begin{figure}[p]\centering
\includegraphics[width=.49\linewidth]{../dumps/sac_jepa_v3_10k_corrected/plots/latent_umap_sac_000010000_seed100001.png}
\includegraphics[width=.49\linewidth]{../dumps/sac_jepa_v3_10k_corrected/plots/dynamics_history_sac_000010000_seed100001.png}
\caption{V3 endpoint held-out terrain/UMAP trajectory (left) and full physical dynamics history (right).}
\end{figure}

\begin{figure}[p]\centering
\includegraphics[width=.49\linewidth]{../dumps/sac_jepa_v3_10k_corrected/plots/encoder_cka_matrix_seed100001.png}
\includegraphics[width=.49\linewidth]{../dumps/sac_jepa_v3_10k_corrected/plots/pairwise_encoder_umap_5x5_seed100001.png}
\caption{V3 encoder evolution across the five checkpoints: CKA (left) and pairwise joint UMAP trajectory matrix (right).}
\end{figure}

\clearpage
\section*{V3 completion: matched 50k objective test (September 21, 2026)}
The preceding 10k V3 result is retained as a source-verified integration
check, but it is not a fair comparison with the 50k V2 pair.  This appended
run uses the same 50,000-transition budget, five checkpoints (10k through
50k), RTX 3090 analysis flush, terrain split, power cap, and SAC settings as
the V2 control.  The only changed learning component is the boundary-safe
$K=8$ action-sequence objective with the inverse-dynamics term described
above.  In particular, time-limit truncations are sequence boundaries, while
they remain bootstrap-able for Bellman learning.

\begin{center}\small
\begin{tabular}{lccc}
\toprule
Run & Transitions & held-out mean speed & final action diagnostic \\
\midrule
SAC-only V2 control & 50k & $2.952\pm0.030$ & --- \\
SAC+JEPA V2, $K=1$ cosine & 50k & $2.954\pm0.049$ & late ratio near 1 \\
SAC+JEPA V3, $K=8$ + inverse & 50k & $2.928\pm0.032$ & ratio $0.684$ \\
\bottomrule
\end{tabular}\end{center}

V3 finishes $0.024$ below SAC-only (0.8\%) and $0.027$ below the V2 cosine
arm (0.9\%).  This single matched run therefore does not show a speed benefit
from the repaired auxiliary task.  It does establish the intended
action-conditioning evidence at the endpoint: real action sequences have
cosine loss $0.0306$ versus $0.0448$ for block-shuffled sequences (ratio
$0.684$), while inverse-action MSE is $0.151$.  The ratio nevertheless rises
close to one late in training before falling again; it is not persistently in
the desired band.  The honest conclusion is that $K=8$ plus inverse dynamics
repairs the previously trivial one-step diagnostic, but needs a multi-seed,
gradient-balanced coefficient sweep before it can be claimed to help control.

\begin{figure}[p]\centering
\includegraphics[width=.90\linewidth]{../dumps/sac_jepa_v3_50k/plots/v3_action_conditioning_diagnostics.png}
\caption{Matched 50k V3 action-conditioning diagnostics.  Top: held-out
speed. Middle: real $K=8$ action sequences versus block-shuffled sequences.
Bottom: their ratio and endpoint inverse-dynamics MSE.}
\end{figure}

\begin{figure}[p]\centering
\includegraphics[width=.49\linewidth]{../dumps/sac_jepa_v3_50k/plots/latent_umap_sac_000050000_seed100001.png}
\includegraphics[width=.49\linewidth]{../dumps/sac_jepa_v3_50k/plots/dynamics_history_sac_000050000_seed100001.png}
\caption{Matched V3 50k endpoint on held-out seed 100001.  Left: periodic
potential, deterministic trajectory, and encoder UMAP. Right: speed, latent
velocity, actions, and signed acceleration components.}
\end{figure}

\begin{figure}[p]\centering
\includegraphics[width=.49\linewidth]{../dumps/sac_jepa_v3_50k/plots/encoder_cka_matrix_seed100001.png}
\includegraphics[width=.49\linewidth]{../dumps/sac_jepa_v3_50k/plots/pairwise_encoder_umap_5x5_seed100001.png}
\caption{Matched V3 representation evolution across checkpoints 10k--50k:
linear CKA (left) and pairwise joint UMAP trajectories on one shared held-out
probe (right).}
\end{figure}

\clearpage
\section*{Retained PPO+JEPA 10k evidence}
For completeness, this report now includes the existing PPO+JEPA 3090 run,
which was omitted from the SAC-focused V3 report.  It uses the earlier
one-step MSE objective
\[
 \|q_\psi([f_\theta(o_t),a_t])-\operatorname{sg}[\bar f(o_{t+1})]\|_2^2
\]
with coefficient $0.1$, so it is not directly comparable to the V3 $K=8$
cosine-plus-inverse objective.  Its recorded 10k held-out result is
$0.676\pm0.105$, compared with $0.818\pm0.116$ for the PPO-only reference.
Thus the earlier PPO auxiliary branch was active but did not improve speed in
that short single-run comparison.  Its retained artifacts include plots but
not the original checkpoints or metrics CSV; the numbers are preserved here
as historical run evidence rather than recomputed claims.

\begin{figure}[p]\centering
\includegraphics[width=.49\linewidth]{../dumps/ppo_jepa_10k/plots/learning_curve.png}
\includegraphics[width=.49\linewidth]{../dumps/ppo_jepa_10k/plots/episode_diagnostics_ppo_000010000_seed100001.png}
\caption{Retained PPO+JEPA 10k learning curve (left) and complete deterministic
episode diagnostic with terrain overlay, controls, and acceleration traces
(right).}
\end{figure}

\begin{figure}[p]\centering
\includegraphics[width=.36\linewidth]{../dumps/ppo_jepa_10k/plots/latent_umap_ppo_000010000_seed100001.png}
\includegraphics[width=.36\linewidth]{../dumps/ppo_jepa_10k/plots/encoder_cka_matrix_seed100001.png}
\par\vspace{0.35em}
\includegraphics[width=.72\linewidth]{../dumps/ppo_jepa_10k/plots/pairwise_encoder_umap_5x5_seed100001.png}
\caption{Retained PPO+JEPA inference and representation artifacts: endpoint
terrain/UMAP correspondence (upper left), CKA across saved checkpoints (upper
right), and pairwise UMAP trajectory matrix (bottom).}
\end{figure}

\clearpage
\section*{Matched PPO-only versus PPO+JEPA at 50k (September 21, 2026)}
This section replaces neither the retained 10k PPO+JEPA evidence nor any SAC
result.  It appends the previously missing fair PPO comparison: two new,
source-verified RTX 3090 runs with identical centered-Gaussian terrain,
random tiled starts, $P_{\max}=6$, 16 environments, seed 123, held-out seeds
100001--100004, 50,000 training transitions, evaluations every 2,048
transitions, and checkpoints at 10k--50k.  PPO+JEPA alone enables the
one-step cosine action-conditioned auxiliary loss with coefficient $0.1$ and
EMA update fraction $0.01$; it is the current PPO branch, not the $K=8$ SAC
V3 objective.

\begin{center}\small
\begin{tabular}{lccc}
\toprule
Run & held-out speed at 50k & JEPA loss at 50k & difference \\
\midrule
PPO-only & $1.658\pm0.268$ & --- & --- \\
PPO+JEPA (cosine, $K=1$) & $1.717\pm0.219$ & $0.00439$ & $+0.058$ ($+3.5\%$) \\
\bottomrule
\end{tabular}\end{center}

The auxiliary PPO run is directionally better in this one matched run.  The
learning curve is also above the control for most of training after the first
few evaluations.  This is evidence that the current PPO+JEPA branch can help
velocity under this protocol, but the held-out standard deviations overlap and
one seed is not a performance claim.  In contrast, the matched SAC runs reach
about $2.95$ speed and remain the strongest controller architecture tested so
far.  The next necessary experiment is a multi-seed PPO replication, followed
by porting the boundary-safe $K=8$ objective to PPO rather than transferring
the SAC V3 conclusion by analogy.

\begin{figure}[p]\centering
\includegraphics[width=.90\linewidth]{../output/pdf/ppo_vs_ppo_jepa_50k_v3.png}
\caption{Matched 50k PPO learning curves.  Shading is the standard deviation
over the four held-out terrain seeds at each deterministic evaluation.}
\end{figure}

\begin{figure}[p]\centering
\includegraphics[width=.49\linewidth]{../dumps/ppo_50k_control_v3/plots/latent_umap_ppo_000050000_seed100001.png}
\includegraphics[width=.49\linewidth]{../dumps/ppo_jepa_50k_v3/plots/latent_umap_ppo_000050000_seed100001.png}
\caption{Endpoint terrain-overlay trajectory and UMAP for matched 50k PPO-only
(left) and PPO+JEPA (right) deterministic inference on held-out seed 100001.}
\end{figure}

\begin{figure}[p]\centering
\includegraphics[width=.49\linewidth]{../dumps/ppo_50k_control_v3/plots/dynamics_history_ppo_000050000_seed100001.png}
\includegraphics[width=.49\linewidth]{../dumps/ppo_jepa_50k_v3/plots/dynamics_history_ppo_000050000_seed100001.png}
\caption{Matched 50k velocity, control, and acceleration histories for PPO-only
(left) and PPO+JEPA (right).}
\end{figure}

\begin{figure}[p]\centering
\includegraphics[width=.49\linewidth]{../dumps/ppo_50k_control_v3/plots/encoder_cka_matrix_seed100001.png}
\includegraphics[width=.49\linewidth]{../dumps/ppo_jepa_50k_v3/plots/encoder_cka_matrix_seed100001.png}
\includegraphics[width=.49\linewidth]{../dumps/ppo_50k_control_v3/plots/pairwise_encoder_umap_5x5_seed100001.png}
\includegraphics[width=.49\linewidth]{../dumps/ppo_jepa_50k_v3/plots/pairwise_encoder_umap_5x5_seed100001.png}
\caption{Representation evolution across the five matched PPO checkpoints:
CKA top row and pairwise UMAP trajectory matrices bottom row; PPO-only left,
PPO+JEPA right.}
\end{figure}

\clearpage
\section*{Spectral-slope transfer for SAC versus SAC+JEPA (September 21, 2026)}
This appended experiment tests whether the SAC comparison depends on the
pink-noise spectral exponent $\alpha$.  The original matched 50k SAC and
SAC+JEPA V2-cosine checkpoints were trained at $\alpha=3$.  Each endpoint was
then used only as an initialization for a fresh 20,000-transition adaptation
run at $\alpha=2$ (more high-frequency terrain) and $\alpha=4$ (more
low-frequency terrain).  Replay buffers, optimizers, entropy temperature, and
live environments were deliberately fresh after the terrain shift; this is
fine-tuning transfer, not an exact training resume.  Both arms retain the
same centered Gaussian-squared potential, random tiled starts, power cap,
seed, and four held-out terrain seeds as the original SAC comparison.

\begin{center}\small
\begin{tabular}{lccc}
\toprule
Spectral exponent & SAC-only speed & SAC+JEPA speed & SAC+JEPA minus SAC \\
\midrule
$\alpha=2$, after 20k transfer & $2.925\pm0.024$ & $2.876\pm0.030$ & $-0.049$ \\
$\alpha=3$, original 50k endpoint & $2.952\pm0.030$ & $2.954\pm0.049$ & $+0.003$ \\
$\alpha=4$, after 20k transfer & $2.841\pm0.119$ & $2.878\pm0.085$ & $+0.037$ \\
\bottomrule
\end{tabular}\end{center}

The overlaid trajectories show fast recovery in every transfer arm.  At
$\alpha=2$, SAC-only ends modestly higher; at $\alpha=4$, SAC+JEPA ends
modestly higher.  These endpoint differences are smaller than, or comparable
to, the across-held-out-terrain variability, and each condition is a single
training seed.  Consequently the correct conclusion is robustness rather
than a JEPA advantage: changing spectral slope did not expose a consistent
SAC+JEPA benefit or regression.  The matched $\alpha=3$ result remains the
primary in-distribution comparison, while this section establishes that both
controllers adapt back toward the $\sim2.95$ speed regime after either slope
shift.

\begin{figure}[p]\centering
\includegraphics[width=.96\linewidth]{../dumps/plots/sac_sac_jepa_spectral_transfer_alpha_2_4.png}
\caption{SAC-only and SAC+JEPA transfer curves after changing the pink-noise
spectral exponent.  Solid lines and shading are held-out mean speed and
standard deviation over four terrain seeds during 20k-transition adaptation.
Dashed lines are the original $\alpha=3$ 50k endpoints, included as a
reference rather than as a matched transfer evaluation.}
\end{figure}

\clearpage
\section*{Controlled start-location replication (September 22, 2026)}
This section appends a stronger check of the preceding spectral-transfer
comparison.  For every $\alpha\in\{2,3,4\}$ and each controller, the
deterministic final checkpoint was evaluated on the same 16 uniformly sampled
locations in each of four held-out terrain tiles (64 episodes per controller).
Within a terrain tile, the reset seed, initial heading, and initial speed were
held fixed, so the only varied initial condition was position.  This protocol
therefore isolates sensitivity to where the agent begins rather than mixing
location variation with a different launch velocity.  The $\alpha=2$ and
$\alpha=4$ controllers were initialized from their $\alpha=3$ 50k weights
and then fine-tuned for 20k transitions with \emph{no encoder freezing};
the actor encoder, critics, and target critics all remained trainable, while
the optimizer and replay state were newly initialized.

\begin{center}\small
\begin{tabular}{lccc}
\toprule
Spectral exponent & SAC-only speed & SAC+JEPA speed & SAC+JEPA minus SAC \\
\midrule
$\alpha=2$ & $2.918\pm0.048$ & $2.882\pm0.040$ & $-0.036$ \\
$\alpha=3$ & $2.991\pm0.072$ & $2.995\pm0.075$ & $+0.004$ \\
$\alpha=4$ & $3.061\pm0.139$ & $3.027\pm0.140$ & $-0.034$ \\
\bottomrule
\end{tabular}\end{center}

This replication changes the practical verdict.  Across common start
locations, SAC-only is higher at both shifted slopes, while the
in-distribution pair is effectively tied.  The $\alpha=2$ and $\alpha=4$
differences are modest relative to episode variation, so they do not prove a
large regression, but neither do they support better JEPA generalization.
The appropriate conclusion is: the current SAC+JEPA auxiliary branch matches
SAC near the training slope but has not demonstrated a robust out-of-slope or
out-of-location benefit.

\begin{figure}[p]\centering
\includegraphics[width=.93\linewidth]{../dumps/spectral_random_start_evaluation/random_start_speed_distribution.png}
\caption{Distribution of deterministic episode mean speeds across 64 shared
start-location evaluations for each spectral exponent and controller.  Green
triangles are means; boxes summarize location and terrain variation.}
\end{figure}

\begin{figure}[p]\centering
\includegraphics[height=.93\textheight]{../dumps/spectral_random_start_evaluation/random_start_trajectory_panels.png}
\caption{Four representative shared-location deterministic trajectories per
controller, overlaid on the corresponding periodic potential.  Each adjacent
SAC/SAC+JEPA pair uses identical first four starts on the same held-out tile.}
\end{figure}

\begin{figure}[p]\centering
\includegraphics[width=.96\linewidth]{../dumps/plots/sac_sac_jepa_spectral_loss_curves.png}
\caption{All logged SAC loss and temperature curves for the original
$\alpha=3$ runs and the $\alpha=2,4$ fine-tuning transfers.  The auxiliary
distance is zero by construction in SAC-only arms.}
\end{figure}

\clearpage
\section*{Frozen-encoder spectral transfer and encoder probes (September 25, 2026)}
This appended comparison freezes the actor visual encoder and both independent
critic visual encoders at the start of transfer.  Only the actor head, Q heads,
temperature, and (where enabled) JEPA prediction heads adapt.  Thus the curves
ask whether the learned representations are useful after a terrain-spectrum
shift, rather than allowing the visual systems to relearn the new terrain.
Vanilla SAC and the original cosine JEPA with a variance-floor regularizer were
run for 50k transfer transitions; the four new regularizer arms were screening
runs of 20k transitions, so their curves stop earlier and are not a matched
50k ranking.

\begin{center}\small
\begin{tabular}{lrrr}
\toprule
Regularizer & $\alpha=1.5$ & $\alpha=2$ & $\alpha=4$ \\
\midrule
VISReg (20k) & 2.549 & 2.198 & 3.066 \\
KerJEPA Gaussian MMD (20k) & 2.762 & 2.594 & 2.902 \\
SUSReg (20k) & 2.973 & 2.971 & 3.004 \\
Spherical MMD (20k) & 3.000 & 2.908 & 2.995 \\
\bottomrule
\end{tabular}\end{center}

The regularizer screening results are mixed: SUSReg is the most consistently
strong at $\alpha=1.5$ and $\alpha=2$ within its 20k budget, while spherical
MMD is strongest at $\alpha=1.5$.  These are single-seed outcomes and should
be treated as candidates for matched 50k replications, not as a conclusive
replacement for the original variance regularizer.

\begin{figure}[p]\centering
\includegraphics[width=.98\linewidth]{../output/pdf/jepa_regularizer_transfer_curves.png}
\caption{Dense held-out speed curves after frozen-encoder transfer to three
spectral exponents.  Shading is the logged held-out standard deviation.  Blue
and orange continue to 50k; the remaining four regularizer screens end at
20k, which is why their right endpoints must not be read as matched 50k
comparisons.}
\end{figure}

\begin{figure}[p]\centering
\includegraphics[width=.95\linewidth]{../dumps/sac_frozen_alpha2_transfer/analysis/actor_q_representation_panel_sac_000020000_seed100001.png}
\caption{Frozen vanilla-SAC transfer at $\alpha=2$: one deterministic held-out
rollout over the potential landscape (upper left), followed by PCA and UMAP of
the actor, Q1, and Q2 encoders.  Every encoder receives the same observation
at each color-coded rollout step; the panels therefore show distinct learned
representations rather than differences caused by separate trajectories.}
\end{figure}

\begin{figure}[p]\centering
\includegraphics[width=.95\linewidth]{../dumps/sac_jepa_actorcritic_frozen_alpha2_transfer/analysis/actor_q_representation_panel_sac_000020000_seed100001.png}
\caption{Frozen actor-and-critic JEPA transfer at $\alpha=2$, evaluated on
the identical held-out terrain seed and deterministic rollout protocol as the
preceding vanilla panel.  Actor, Q1, and Q2 are probed separately before each
action; this verifies that the report distinguishes the policy representation
from both value representations.}
\end{figure}

\end{document}
